% v15: diamond exclusion. Entire proof is analytic and has no denominator bound.
\section{Diamond 型高分母 support の排除}
\label{dg15:section}

高分母 support が四頂点上の diamond であり，第五の core は全低分母である場合を考える．
この場合には，四本の低分母辺に沿う対向係数の相殺を加えることで，
二つの二価 core の低低 entry を同期させることができる．
以下では分母の数値上限も，三高分母局所模型の分類も用いない．

\begin{lemma}[全低分母 core における総和恒等式]
\label{dg15:hub-balance}
五つの core を $0,1,2,3,4$ と標識し，$q_{i4}=1$ $(0\leq i\leq3)$ とする．
各 core の整数行列は行和と pair products を満たし，四辺 $i4$ で対向 $D$ が相殺するとする．
$I=\{0,1,2,3\}$ とし，$i\ne j$ に対して
\[
 x_{ij}=p^i_{j4},\qquad r_{ij}=p^i_{4j},\qquad V_{ij}=p^4_{ij}
\]
と置く．このとき
\begin{equation}\label{dg15:balance-total}
 \sum_{i\ne j\in I}x_{ij}=\sum_{i\ne j\in I}r_{ij}=-4.
\end{equation}
さらに，高分母辺 $ij$ の puncture $4$ で整数 clutching が成立すれば，
\begin{equation}\label{dg15:low-defects}
 \sum_{\substack{i<j\in I\\q_{ij}=1}}
 (x_{ij}+x_{ji}-r_{ij}-r_{ji})=0.
\end{equation}
\end{lemma}

\begin{proof}
各 core の第 $4$ 行から $\sum_{j\ne i}r_{ij}=-1$ を得る．一方，
\[
 D^i_4=-1-\sum_{j\ne i}x_{ij},\qquad
 D^4_i=-1-\sum_{j\ne i}V_{ji}
\]
なので，相殺を $i$ について加えると
$\sum x_{ij}+\sum V_{ij}=-8$ となる．$V$ の四行の和は $-4$ だから
\eqref{dg15:balance-total} が従う．

$q=q_{ij}>1$ のとき，$x_{ij},x_{ji}$ は $q-1$ の符号付き約数であり，
$q\mid x_{ij}+x_{ji}$ である．従ってその対は
\begin{equation}\label{dg15:low-puncture-options}
 (t,-t),\qquad \{1,q-1\},\qquad\{-1,-(q-1)\}
\end{equation}
のいずれかである．実際，和の絶対値は $2(q-1)<2q$ であり，
同符号の和が $q$ なら，一方の約数が $q-1$，他方が $1$ でなければならない．
両方が真の正約数なら，その和は $q-1$ 以下になるからである．
各場合に
\[
 r_{ij}+r_{ji}=(q-1)/x_{ij}+(q-1)/x_{ji}=x_{ij}+x_{ji}
\]
が成立する．高分母辺の寄与を総和から除けば \eqref{dg15:low-defects} を得る．
\end{proof}

\begin{theorem}[Diamond の排除]
\label{dg15:no-diamond}
全平面五次の標準 $K_4$ cones と $K_5$ branch graph に対する，正の整数局所 charts，
marked clutching，対向 $D$ の相殺，node 非負性を満たすデータについて，
高分母 support が
\[
 H=\{01,02,03,12,13\}
\]
となることはない．
この排除には分母の上限を要しない．
\end{theorem}

\begin{proof}
二高分母の標準形と補題~\ref{ml12:dichotomy} の二分を用い，
\begin{equation}\label{dg15:canonical}
 P=\begin{pmatrix}
 0&-A&x&-s\\
 -B&0&-y&z\\
 r&-r-1&0&0\\
 -S&T&-L&0
 \end{pmatrix},\qquad
 \begin{gathered}
 a=sS+1,\quad b=zT+1,\\
 x=(a-1)/r,\quad y=(b-1)/(r+1),\\
 A=a+x-s,\quad B=b-y+z,\quad AB=ab-1,\\
 L=T-S+1
 \end{gathered}
\end{equation}
と書く．全ての $r,s,z,S,T,x,y$ は正整数，$S\geq2$ である．
負の枝は $r>S=T$ で $L=1$，正の枝は $S>r$ で
\begin{equation}\label{dg15:positive-range}
 S<r(L-1)
\end{equation}
を満たす．従って $L\in\{1,3,4,\ldots\}$ である．
以下の primed 記号は第二の標準形を表す．

\paragraph{二つの低低 entry の同期．}
\[
 u_2=p^2_{34},\quad v_2=p^2_{43},\qquad
 u_3=p^3_{24},\quad v_3=p^3_{42}
\]
と置く．二高分母の標準形により，各対 $(u_i,v_i)$ は $(0,-L_i)$ または
$(-L_i,0)$ であり，$L_i>0$ である．補題~\ref{dg15:hub-balance} の総和には
低分母辺 $23$ だけが残るので，$u_2+u_3=v_2+v_3$ となる．
従って $L_2=L_3=:L$ であり，向きは反対である．$2,3$ を交換して
\[
 (u_2,v_2)=(-L,0),\qquad (u_3,v_3)=(0,-L)
\]
とする．node $23;4$ の数は $1+L$ だから $L\leq4$，従って
\begin{equation}\label{dg15:three-L}
 L\in\{1,3,4\}.
\end{equation}

core $2$ の第 $3$ 行と core $3$ の第 $2$ 行には，高 ray に向かう正の entry が
それぞれ一つだけある．もしその行先が同じなら，対応する node $23;0$ または $23;1$ の数が負になる．
従って $0,1$ を交換すれば，core $2$ の標準形は $(B,A,D,C)=(0,1,3,4)$，
core $3$ は $(B,A,C,D)=(0,1,2,4)$ とできる．
辺 $23$ の三つの node 数は
\begin{equation}\label{dg15:n23}
 (n_{23;0},n_{23;1},n_{23;4})
 =(r'+2-T,\ 1+S-r',\ 1+L).
\end{equation}
従って
\begin{equation}\label{dg15:rprime-range}
 S+L-3\leq r'\leq S+1.
\end{equation}

\paragraph{共通低 puncture の例外の全列挙．}
中心 core $0,1$ の第 $4$ 行で，共有高 ray に向かう entry を
\[
 R_0=p^0_{41},\qquad R_1=p^1_{40}
\]
と書く．$q_{01}=Q>1$ なので，これらは非零であり，
\eqref{dg15:low-puncture-options} と同じ三択を満たす．特に，異符号なら和は零である．

四辺 $02,12,03,13$ の puncture $4$ で \eqref{dg15:low-puncture-options} の
反対符号の場合を選ぶと，中心側の第 $4$ 行は
\[
 p^0_{42}=r+1,\quad p^1_{42}=-r,\qquad
 p^0_{43}=-T',\quad p^1_{43}=S'.
\]
その行和から
\begin{equation}\label{dg15:central-baseline}
 R_0=T'-r-2,\qquad R_1=r-S'-1,\qquad R_0+R_1=L-4
\end{equation}
を得る．この選択と異なる場合を以下では例外と呼ぶ．全ての例外は次の表で尽くされる．
表の outgoing entry は，二価側の共有高 ray の行から puncture $4$ に向かう成分を指す．
\[
\begin{array}{c|c|c|c}
 \text{辺}&\text{二価側の outgoing entry}&\text{例外の必要条件}&
 \text{中心の第4行の変化}\\\hline
 02&-y&y=1&p^0_{42}:\ r+1\longmapsto-1\\
 12&x&r=1\ \text{または}\ x=1&p^1_{42}:\ -r\longmapsto -r+a\\
 03&z'&z'=1&p^0_{43}:\ -T'\longmapsto1\\
 13&-s'&s'=1&p^1_{43}:\ S'\longmapsto-1
\end{array}
\]
例えば $y=(b-1)/(r+1)<b-1$ だから，負の同符号対は $y=1$ の場合だけである．
他の三行も \eqref{dg15:low-puncture-options} と $S',T'\geq2$ から従う．
例外による $(R_0,R_1)$ の変化は，順に $(b,0),(0,-a),(-b',0),(0,a')$ である．

ここで $x=1$ は常に不可能である．正の枝では $x=sS/r>s\geq1$ である．
負の枝では $H=r-S$，$u=s-x$，$v=z-y$ と置くと，正整数 $u,v$ は
既出の式~\eqref{ml12:E-equation}，すなわち
\begin{equation}\label{dg15:negative-identity}
 Euv=v-u+1,\qquad E=1-\frac{S^2}{H(H+1)}
\end{equation}
を満たす．$u=1$ は $E=1$ を要求するので $u\geq2$ である．
$u(1+Ev)=v+1$ より $E>-1$，また $E>0$ なら
$Eu=1-(u-1)/v<1$ より $E<1/2$ である．従って
\begin{equation}\label{dg15:negative-range}
 \frac{S^2}{2}<H(H+1)<2S^2.
\end{equation}
なお $S=1$ は $H(H+1)\geq2$ と右側の厳密不等式に反するので，負の枝でも $S\geq2$ である．
$x=1$ なら $r=sS$，$s=u+1\geq3$，従って $H\geq2S$ となり，
右側の不等式に反する．これで上表の全例外を吟味する準備ができた．

\paragraph{$L=1$ の排除．}
両方とも負の枝であり，$r>S\geq2$，$r'>S'\geq2$ である．
従って $r=1$ は不可能であり，$s'\geq2$，$z'\geq2$ なので，残る例外は $y=1$ だけである．

例外がなければ $R_0+R_1=-3$ なので，\eqref{dg15:low-puncture-options} から
$Q=3$，$R_0,R_1<0$ となる．特に $r\leq S'$ である．しかし
\[
 S<r\leq S'<r'\leq S+1
\]
は整数について不可能である．

$y=1$ の例外を選ぶ場合は $b=r+2$ なので
\[
 R_0=S'\geq2,\qquad R_0+R_1=r-1>0.
\]
同符号正の対は $\{1,Q-1\}$ だから，$R_1=1$，$S'=r-2$ である．
$r'\geq S'+1=r-1$ と \eqref{dg15:rprime-range} から $S\geq r-2$，
従って $H=r-S\in\{1,2\}$ となる．さらに $y=1$ は
\[
 (z-1)S=H+1
\]
を与える．$S\geq2$ より $(H,S)=(1,2),(2,3)$ の二つだけである．
前者は \eqref{dg15:negative-range} に反する．後者では
$(r,S,S',r')=(5,3,3,4)$ となり，primed 模型の $H'=1,S'=3$ が
同じ不等式に反する．

\paragraph{$L=3$ の排除．}
両方とも正の枝であり，$r<S<2r$，$r'<S'<2r'$ である．
従って $r=1$ は不可能である．また $y=zT/(r+1)>z\geq1$ なので $y=1$ は不可能である．
$s'=1$ なら $x'=S'/r'\in(1,2)$ となる．
$z'=1$ なら
\[
 y'=\frac{S'+2}{r'+1}\in(1,2)
\]
となる．いずれも整数性に反する．従って例外は一つもなく，
$R_0+R_1=-1$ が $Q>1$ による可除性に反する．

\paragraph{$L=4$ の例外の排除．}
この場合 $r'=S+1\geq3$ である．引き続き $y=1$ は不可能である．
$s'=1$ なら $1<S'/r'<3$ と整数性から $S'=2r'$，$T'=2r'+3$ となる．
$y'$ の整数性より $z'=h(r'+1)$，$h\geq1$ と書ける．
$A'B'=a'b'-1$ へ代入すると
\[
 h(r'+1)=2
\]
となり，$r'=1$ を要求するため不可能である．

$z'=1$ なら $1<T'/(r'+1)<3$ なので
$T'=2(r'+1)$，$S'=2r'-1$ となる．$x'$ の整数性より $s'=h r'$ と書け，
同じ pair equation は
\[
 h(r'-2)=0
\]
となる．従って $r'=2$ を要求し，これも不可能である．
これら二つの代入恒等式は任意の正整数に対する恒等式であり，有界探索を用いていない．

残る例外は $r=1$ の場合の辺 $12$ だけである．これを選ぶと
$S=2$，$T=5$，$r'=3$ であり，
\[
 R_0=S'>0,\qquad R_1=-S'-a<0.
\]
異符号でありながら和が $-a\ne0$ となるので，この例外も不可能である．

\paragraph{$L=4$ の反対符号の場合の排除．}
残った場合は \eqref{dg15:central-baseline} そのものであり，
\[
 R_0=R=S'-r+1,\qquad R_1=-R.
\]
$S\geq r+1$，$r'=S+1$，$S'\geq r'+1$ より $R\geq4$ である．
全低 core の行列を引き続き $V$ と書く．node $04;1$ の非負性は
$V_{01}\leq1-R<0$ を与える．$V_{01}V_{10}=0$ より $V_{10}=0$ である．
従って node $14;0$ の数は $1+R$ であり，辺 $14$ の node 総数五より $R\leq4$ となる．
結局
\[
 R=4,\qquad S=r+1,\qquad r'=r+2,\qquad S'=r+3.
\]
node $14;0$ が五を使い尽くすので $n_{14;2}=0$ であり，
$p^1_{42}=-r$ より $V_{12}=r+1>0$ となる．
今度は $V_{12}V_{21}=0$ より $V_{21}=0$ だから，
\[
 n_{24;1}=1-p^2_{41}-V_{21}=1-r\geq0.
\]
従って $r=1$，すなわち primed 模型は $(r',S',T')=(3,4,7)$ でなければならない．

この模型では $s'=3h$，$z'=4k$，$h,k\geq1$ と書け，
\[
 a'=12h+1,\quad b'=28k+1,\quad A'=13h+1,\quad B'=25k+1.
\]
従って pair equation は
\[
 k(11h+3)=h+1
\]
となる．左辺は $11h+3>h+1$ 以上なので不可能である．
以上で \eqref{dg15:three-L} の全場合を排除した．
\end{proof}

\begin{remark}[使用した大域条件の範囲]
\label{dg15:scope}
上の証明で $D$ の相殺を使うのは全低 core に接する四辺 $i4$ の総和だけである．
二つの三高分母中心の正符号 graph の分類や，中心側の正 weights の存在判定は用いていない．
二価 core の局所標準形，低 puncture の整数 clutching，行和，pair products，
指定した低辺の node 非負性だけを組み合わせている．
従って，三高分母の matching 定理や全高 $D$ の符号に依存しない．
\end{remark}
